\chapter{Crise, colapso e poder}
\noindent\textit{Vinheta.} Numa terça-feira comum, um ataque de \emph{ransomware} paralisa uma operadora de logística que atende supermercados, farmácias e postos. Caminhões param, sistemas de estoque caem, prazos explodem. Em poucas horas, não é ``uma empresa com problema'': é a cidade que fica mais \textbf{frágil}. A narrativa oficial pede calma; nos bastidores, decisões emergenciais definem quem recebe primeiro. \textbf{Poder} flui para quem controla as chaves.

\section{Mapa bíblico}
\begin{itemize}[leftmargin=*,noitemsep]
  \item \textbf{Poder e serviço} (Mc~10:42--45): governar como quem serve, não dominar.
  \item \textbf{Cidades que caem} (Ap~18): cadeias econômicas desmoronam ``\emph{numa hora}''.
  \item \textbf{Vigilância e sobriedade} (1~Pe~5:8; Mt~24): atenção e perseverança, sem pânico.
  \item \textbf{Justiça e verdade} (Pv~21:3; Ef~4:25): transparência e retidão em tempos de confusão.
\end{itemize}

\section{Três motores de crise}
\begin{itemize}[leftmargin=*,noitemsep]
  \item \textbf{Ciberataques}: \emph{ransomware}, DDoS, comprometimento de cadeia de suprimentos de software.
  \item \textbf{Choques físicos/logísticos}: clima extremo, greves, bloqueios, acidentes em portos/rodovias.
  \item \textbf{Crises de confiança/narrativa}: boatos, \emph{deepfakes}, decisões opacas que geram corrida a mercados/bancos.
\end{itemize}

\section{Lei marcial informacional (o que é e como lidar)}
Chamamos de \textbf{lei marcial informacional} o conjunto de \emph{restrições emergenciais} sobre comunicação e acesso a plataformas durante crises (para conter pânico, golpes ou desordem). Pode incluir limitação de \emph{upload}, derrubada de conteúdos, prioridade a canais oficiais e suspensão temporária de certos recursos.
\begin{itemize}[leftmargin=*,noitemsep]
  \item \textbf{Risco}: decisões opacas, abuso de poder, silenciamento de alertas legítimos.
  \item \textbf{Resposta cristã}: buscar o bem comum sem idolatrar a ordem; falar verdade com prudência; garantir \emph{canais alternativos} para cuidado.
\end{itemize}

\begin{nicbox}
\textbf{NIC aplicado}\\
\textbf{Narrativa}: quem define o que é ``oficial''? há prestação de contas?\\
\textbf{Identidade}: quem ganha o rótulo de ``confiável'' ou ``problemático''?\\
\textbf{Controle}: quem recebe \emph{prioridade} de recursos (combustível, tráfego, reposição)? por quais critérios?
\end{nicbox}

\section{Quadro de sinais precoces (o que acompanhar \emph{antes} de faltar no mercado)}
\begin{checklist}[title={Indicadores práticos para leigos}]
\begin{itemize}[leftmargin=*,noitemsep]
  \item \textbf{Energia}: comunicados de racionamento, avisos de manutenção extensa, quedas em múltiplos bairros.
  \item \textbf{Telecom}: falhas simultâneas de mais de uma operadora; latência anormal; apps bancários instáveis.
  \item \textbf{Logística}: notícias de porto/rodovia fechados; filas em postos; relatos de falta de insumos (gás, diesel).
  \item \textbf{Comércio}: limite de compras por item; “apenas dinheiro”; prateleiras recorrentes vazias no \emph{básico} (arroz, leite, ovos).
  \item \textbf{Saúde}: aumento de espera em UPA/UBS; falta de medicamentos comuns; suspensão de exames eletivos.
  \item \textbf{Narrativa}: mudança de tom oficial para ``\emph{pequenas restrições temporárias}'' sem prazo claro.
\end{itemize}
\end{checklist}

\section{Cenários críticos e contramedidas}
\subsection*{1) Ciberataque a pagamentos/ERP}
\begin{itemize}[leftmargin=*,noitemsep]
  \item \textbf{Efeitos}: cartão/TEF fora; notas fiscais indisponíveis; conciliação quebrada.
  \item \textbf{Contramedidas}: dinheiro trocado; recibo manual; compras pequenas; priorize comércios que operam sem sistema.
\end{itemize}

\subsection*{2) Paralisação logística (portos/rodovias)}
\begin{itemize}[leftmargin=*,noitemsep]
  \item \textbf{Efeitos}: desabastecimento pontual (combustível, hortifrúti); alta de preços.
  \item \textbf{Contramedidas}: reduzir deslocamentos; compras fracionadas; feiras locais; cozinhas compartilhadas.
\end{itemize}

\subsection*{3) Corte de energia em larga escala}
\begin{itemize}[leftmargin=*,noitemsep]
  \item \textbf{Efeitos}: cadeia fria comprometida; bombas d'água paradas; telecom instável.
  \item \textbf{Contramedidas}: priorize perecíveis; água estocada; \emph{power banks}; vizinhança solidária.
\end{itemize}

\subsection*{4) Lei marcial informacional (restrição de canais)}
\begin{itemize}[leftmargin=*,noitemsep]
  \item \textbf{Efeitos}: difícil verificar boatos; canais oficiais saturados; confusão.
  \item \textbf{Contramedidas}: pontos de encontro; líderes de quarteirão; boletins impressos; rádio/comunicação simples.
\end{itemize}

\section{Protocolo 48--72 horas (família e igreja)}
\begin{checklist}[title={Passo a passo objetivo},breakable]
\begin{itemize}[leftmargin=*,noitemsep]
  \item \textbf{Hora 0--6}: confirmar segurança de todos; desligar equipamentos sensíveis; definir ponto de encontro e horário de checagem.
  \item \textbf{Hora 6--12}: água e alimentos essenciais; contactar vizinhos vulneráveis; separar dinheiro trocado; organizar medicações.
  \item \textbf{Hora 12--24}: racionar uso de bateria; cozinhar perecíveis; registrar ocorrências (fotos, anotações) e \textbf{não} espalhar boatos.
  \item \textbf{Hora 24--48}: reavaliar estoque; planejar revezamento de tarefas (cozinhar, buscar água, vigiar); boletim simples para o grupo/igreja.
  \item \textbf{Hora 48--72}: se persistir, ativar redes de apoio (compra coletiva de alimentos, acolhimento de famílias); culto doméstico/comunitário de encorajamento.
\end{itemize}
\end{checklist}

\section{Ferramenta --- Política de crise da igreja (uma página)}
\begin{checklist}
\begin{itemize}[leftmargin=*,noitemsep]
  \item \textbf{Cadastre líderes de quarteirão}: 1 contato por rua/condomínio.
  \item \textbf{Defina ponto físico de referência}: templo ou casa estratégica.
  \item \textbf{Escreva mensagens-tipo}: ``Reconhecemos X; estamos fazendo Y; pedimos Z; horário do próximo boletim.''
  \item \textbf{Critérios de auxílio}: prioridades para idosos, crianças, doentes e trabalhadores essenciais.
  \item \textbf{Transparência}: prestação de contas pública de doações e gastos.
\end{itemize}
\end{checklist}

\section{Postura do discípulo em tempos de colapso}
\begin{itemize}[leftmargin=*,noitemsep]
  \item \textbf{Coração}: esperança firme (cap.~4), sem cinismo.
  \item \textbf{Cabeça}: informação verificada (cap.~3 e 6), sem credulidade.
  \item \textbf{Mãos}: serviço prático (cap.~10), sem egoísmo.
\end{itemize}

\section{Perguntas para grupos pequenos}
\begin{itemize}[leftmargin=*,noitemsep]
  \item Quais \textbf{sinais precoces} você já percebeu na sua cidade nos últimos meses?
  \item Sua família consegue cumprir este \textbf{protocolo 48--72h} hoje? O que falta?
  \item Como sua igreja pode decidir \textbf{critérios de auxílio} justos antes da próxima crise?
\end{itemize}

\bigskip
\noindent\textit{Próximo capítulo: O Anticristo --- paródia da Trindade e o mínimo que todo cristão precisa saber.}
